\color{unidarkgreen}{\subsection{Логика деления стороны СН (ЛД СН)}}
\color{black}

\subsubsection{Общие сведения}

\begin{enumerate}[label=\arabic{section}.\arabic{subsection}.\arabic{subsubsection}.\arabic*,  align=left, labelsep=4pt, leftmargin=0pt, itemindent=57pt]

\begin{figure}[H]
\centering
\includegraphics[width=0.8\textwidth,height=0.8\textheight,keepaspectratio]{\fbpath/022.1791 ЛД СН/img/LD_SN.pdf}
\captionof{figure}{Функциональная схема ЛД стороны СН}\label{pic:LD_SN}
\end{figure}

\item
Ввод в работу логики деления средней стороны трансформатора осуществляется с помощью
уставок \textbf{«ЛД СН от МТЗ-1 СН»} и \textbf{«ЛД СН от МТЗ-2 СН»}. Оперативный ввод/вывод ЛД средней стороны
производится с помощью БУ (блока управления функциями) «Вывод ЛД СН».   

\item
От направленных пусков МТЗ СН в сторону сети «П. МТЗ-1 СН в сеть» и «П. МТЗ-2 СН в сеть»
через выдержки времени \textbf{«Тср МТЗ-1 откл СВ(ШСВ) СН»} и \textbf{«Тср МТЗ-2 откл СВ(ШСВ) СН»} логика
деления формирует сигналы на отключение СВ (ШСВ) средней стороны трансформатора с
возможностью контроля включенного положения выключателя средней стороны.

\item
От направленных пусков МТЗ СН в сторону сети «П. МТЗ-1 СН в сеть» и «П. МТЗ-2 СН в сеть»,
через выдержки времени \textbf{«Тср МТЗ-1 откл В СН»} и \textbf{«Тср МТЗ-2 откл В СН»} логика деления формирует
сигналы на отключение выключателей средней стороны трансформатора.

\small
\begin{longtable}{|>{\centering\arraybackslash}m{5.3cm}|>{\centering\arraybackslash}m{4.2cm}|>{\centering\arraybackslash}m{1cm}|>{\centering\arraybackslash}m{3.3cm}|>{\centering\arraybackslash}m{1.8cm}|}
% Первая страница: полный заголовок
\caption{Перечень уставок ЛД стороны СН\hfill\vspace{-0.5\baselineskip}}\label{desc:tbl1}\\
\hline
\rowcolor{gray!30}
Наименование уставки & Диапазон & Шаг & Значение по умолчанию & Единица измерения \\
\hline
\endfirsthead

% Промежуточные страницы: «Продолжение таблицы»
\caption*{\hspace{3pt}\emph{Продолжение таблицы \ref{desc:tbl1}}\hfill\ }\vspace{-0.5\baselineskip} \\
\hline
\rowcolor{gray!30}
Наименование уставки & Диапазон & Шаг & Значение по умолчанию & Единица измерения \\
\hline
\endhead

%===m>ДТЗ|-
%===m
\end{longtable}
\normalsize	

\end{enumerate}